% Experiment report. Fixed sections: see experiments/TEMPLATE/report.tex and experiments/README.md.
% Build: npm run experiment:report -- 03-front-pull
\documentclass{../polyexp}
\begin{document}
% generated by scripts/experiment-report.ts from experiment.json and models/<run>/ - do not edit
\providecommand{\ExpHeader}{}\renewcommand{\ExpHeader}{%
{\sffamily\footnotesize\color{expgrey}EXPERIMENT 3 \textperiodcentered\ MECHANISM \hfill 2026-11-24}\par
{\LARGE\bfseries\sffamily Pull idle units to the front}\hfill\verdictbox{expgreen}{KEPT}\par\smallskip
{\itshape Does the bot get stronger when idle units drift toward the nearest enemy?}\par\medskip
\begin{tabular}{@{}p{22mm}p{\dimexpr\linewidth-26mm}@{}}
\fact{Change}{A move scores +front per tile it brings the unit closer to the nearest known enemy city or unit. The centre pull remains the fallback when no enemy is known.}
\fact{Knobs}{\knob{front} $\in$ \{0.5, 1, 2, 4, 8\}, chosen 1.13}
\fact{Base}{\model{greedy-full-v1} $\rightarrow$ result \model{greedy-full-front}}
\fact{Method}{Tuning run (100,000 games), then the first strength sweep (200 games per point, final test 400); metric: points share}
\fact{Setups}{4 players, 16$\times$16, drylands, turn cap 150}
\fact{Record}{issue \issue{30}, pull request \issue{62}; status: done}
\end{tabular}\par\medskip
\begin{center}\begin{tikzpicture}[node distance=5mm]
    \node[flowbox] (n0) {{\sffamily\bfseries\scriptsize\color{expgrey}BASE}\\[1pt]\model{full-v1}};
    \node[flowbox,right=5mm of n0] (n1) {{\sffamily\bfseries\scriptsize\color{expgrey}CHANGE}\\[1pt]\knob{front}\\{\scriptsize 0.5, 1, 2, 4, 8}};
    \node[flowbox,right=5mm of n1] (n2) {{\sffamily\bfseries\scriptsize\color{expgrey}MEASURED}\\[1pt]Tuning run, then the first strength sweep\\{\scriptsize 1 setup, metric: points share}};
    \node[flowbox,right=5mm of n2] (n3) {{\sffamily\bfseries\scriptsize\color{expgrey}RESULT}\\[1pt]Points vs base, final test\\{\scriptsize 50.0 $\rightarrow$ 54.8\,\%}\\\textbf{+4.8\pp}};
    \node[flowbox,right=5mm of n3] (n4) {{\sffamily\bfseries\scriptsize\color{expgrey}VERDICT}\\[1pt]\textbf{\color{expgreen}KEPT}\\{\scriptsize new champion:}\\\model{full-front}};
    \draw[flowarrow] (n0) -- (n1); \draw[flowarrow] (n1) -- (n2); \draw[flowarrow] (n2) -- (n3); \draw[flowarrow] (n3) -- (n4);
\end{tikzpicture}\end{center}
\begin{center}\sffamily\small\begin{tabular}{@{}lrrr@{}}\toprule
Measure & Base & New & Change\\\midrule
Points vs base, final test & 50.0\,\% & 54.8\,\% & +4.8\\
Won vs base, final test & 50.0\,\% & 57.8\,\% & +7.8\\
Points share, 6-model panel & 67.7\,\% & 70.3\,\% & +2.6\\
\bottomrule\end{tabular}\end{center}
}
\providecommand{\ExpResults}{}\renewcommand{\ExpResults}{%
\subsection{Tuning run, 100,000 games (the new champion) {\normalfont\small(\model{greedy-full-front})}}
\begin{center}\begin{tikzpicture}\begin{axis}[expplot,height=4.6cm,xlabel={training games (thousands)},ylabel={Points share, panel mean (\%)},enlarge y limits=0.2]
\addplot[dashed,expgrey,line width=0.7pt,forget plot] coordinates {(0.0,68.10) (100.3,68.10)};
\addplot[expblue,mark=*,mark size=1.8pt] coordinates {(0.0,68.10) (12.5,68.10) (25.4,68.10) (37.6,68.10) (50.2,68.10) (62.9,68.10) (75.4,70.30) (87.8,70.30) (100.3,70.30)};
\end{axis}\end{tikzpicture}\\[-2pt]{\footnotesize\color{expgrey}Learning curve of the tuning run (2 kept changes in 258 rounds). Dashed: the start, i.e. the mechanism at its hand-set value.}\end{center}
\begin{center}\sffamily\small\begin{tabular}{@{}lrr@{}}\toprule
Opponent & Won (\%) & Points (\%)\\\midrule
\model{v1} & 94 $\pm$ 2 & 75.7\\
\model{hill-v1} & 91 $\pm$ 3 & 72.3\\
\model{hill-v2} & 95 $\pm$ 2 & 73.7\\
\model{hill-v3} & 91 $\pm$ 3 & 70.3\\
\model{hill-v4} & 87 $\pm$ 3 & 68.6\\
\model{full-v1} & 58 $\pm$ 5 & 54.8\\
\midrule
\textbf{Mean} & 85.9 & 69.2\\
\bottomrule\end{tabular}\\[3pt]{\footnotesize\color{expgrey}Final test of the saved model: 400 unseen games against each panel model.}\end{center}

\subsection{Strength sweep, untrained {\normalfont\small(\model{greedy-full-front-sweep})}}
\begin{center}\sffamily\small\begin{tabular}{@{}lrrrrrr@{}}\toprule
Point & Finish (\%) & Points (\%) & Won (\%) & Conquest (\%) & Turns & vs base: points (\%)\\\midrule
base & -- & 67.73 & 84.3 & -- & -- & 50.0\\
front=0.5 & -- & 68.50 & 86.0 & -- & -- & 53.0\\
front=1 & -- & 68.13 & 85.3 & -- & -- & 53.3\\
front=2 & -- & 68.77 & 86.8 & -- & -- & 52.9\\
front=4 & -- & 68.43 & 86.2 & -- & -- & 52.1\\
front=8 & -- & 68.55 & 86.7 & -- & -- & 53.4\\
\bottomrule\end{tabular}\\[3pt]{\footnotesize\color{expgrey}Panel means; 200 games against each of 6 panel models per point.}\end{center}
\begin{center}\begin{tikzpicture}\begin{axis}[expplot,height=4.6cm,xtick={0,1,2,3,4,5},xticklabels={{base},{0.5},{1},{2},{4},{8}},
  xlabel={\texttt{front}},ylabel={Points share, panel mean (\%)},enlarge x limits=0.12,enlarge y limits=0.25]
\addplot[dashed,expgrey,line width=0.7pt,forget plot] coordinates {(0,67.733) (5,67.733)};
\addplot[only marks,mark=*,mark size=2.6pt,expblue] coordinates { (1,68.500) (2,68.133) (3,68.767) (4,68.433) (5,68.550)};
\addplot[only marks,mark=*,mark size=2.6pt,expgrey] coordinates {(0,67.733)};
\end{axis}\end{tikzpicture}\\[-2pt]{\footnotesize\color{expgrey}Dashed: the base. Grey: base; blue: sweep points.}\end{center}
\begin{center}\sffamily\small\begin{tabular}{@{}lrr@{}}\toprule
Opponent & Won (\%) & Points (\%)\\\midrule
\model{v1} & 96 $\pm$ 2 & 75.5\\
\model{hill-v1} & 94 $\pm$ 2 & 72.2\\
\model{hill-v2} & 95 $\pm$ 2 & 72.3\\
\model{hill-v3} & 91 $\pm$ 3 & 69.1\\
\model{hill-v4} & 88 $\pm$ 3 & 67.0\\
\model{full-v1} & 60 $\pm$ 5 & 54.1\\
\midrule
\textbf{Mean without \model{full-v1}} & 92.5 & 71.2\\
\bottomrule\end{tabular}\\[3pt]{\footnotesize\color{expgrey}Final test of front=2: 400 unseen games against each.}\end{center}
}

\ExpHeader

\section{Question and hypothesis}
A move scores only if it brings the unit closer to its goal. A unit whose goal is far away or unreachable stays where it is, often at the board edge. \textbf{Hypothesis:} a pull toward the nearest enemy puts idle units to work and raises strength.

\section{Mechanism}
Let $d(x)$ be the distance from tile $x$ to the nearest enemy city or unit the player has seen (not at peace). A move from $a$ to $b$ gets $+\,\knob{front}\cdot(d(a)-d(b))$. Default 0. Code: \texttt{enemy\_dist}, \texttt{front\_term} in \texttt{bot\_full.rs}.

\section{Setup}
The tuning run started from \knob{front} $=1$ under the protocol of experiments 1 and 2. After it, the strength sweep was introduced: the base plus the knob at five hand-set values, each against the panel on the same maps as the curve points, then a final test of the best point. The sweep took 1\,min 54\,s; a tuning run takes about 30 minutes.

\section{Results}
\ExpResults

The trained model reaches 70.3\,\% on the panel (base 67.7\,\%, control 66.0\,\%) and beats the base head to head: 54.8\,\% of points, 58\,\% won ($\pm$5). The sweep shows the gain hardly depends on the strength: every value from 0.5 to 8 adds 0.4 to 1.1 points on the panel and about 53\,\% against the base. Training added about 1.5 points more, from two kept changes late in the run that also moved \knob{open.attack}, \knob{greedy.richAt} and \knob{greedy.hit}.

\section{Behaviour}
The behaviour metric planned for this step was not measured; the strength result was judged clear enough on its own. This is a gap in the record. Planned metric: mean distance of units to the nearest enemy.

\section{Interpretation}
The pull mostly gets idle units moving at all, and any strength does that; this is why the response is flat across a factor of 16. The first mechanism of the series that adds something the bot could not do before, instead of re-weighting a decision it already made.

\todo{jakob}{Your reading: does this match what you saw? Anything to add or correct?}

\section{Threats to validity}
\begin{itemize}
\item The champion's extra 1.5 points come from two late kept changes in one run; a different seed might not find them.
\item Under the finish score (experiment 5) a stronger pull trades points for faster conquests, so ``flat'' holds for points share only.
\end{itemize}

\section{Verdict and next step}
\textbf{Kept} (head-to-head rule). New champion: \model{greedy-full-front}. Process change: new mechanisms are measured with a strength sweep; tuning runs happen every three to four kept mechanisms. Next: road-aware moves (experiment 4).

\section{Notes for the paper}
Draft: mechanisms that enable a new behaviour (move when idle) give gains that are insensitive to their weight; a 2-minute sweep finds them as reliably as a 30-minute tuning run. Candidate figure: the sweep plot (flat response) next to the learning curve (170 rounds with nothing kept).

\todo{jakob}{Is this worth a place in the paper? Where?}

\end{document}
